% =============================================================
% THESIS EXTENDED TIKZ STYLES
% File: thesis_tikz_styles.tex
% Purpose: Power-system-specific TikZ styles, pgfplots themes,
%          and macros for the PhD thesis
% =============================================================

%% ── Additional TikZ libraries ────────────────────────────────
\usetikzlibrary{%
  shapes, shapes.geometric, shapes.symbols,
  arrows, arrows.meta,
  calc, positioning, fit,
  decorations.pathmorphing,
  decorations.pathreplacing,
  decorations.markings,
  patterns,
  matrix,
  mindmap,
  trees,
  shadows,
  3d,
  perspective,
}

%% ── Thesis colour palette ────────────────────────────────────
\definecolor{thesisblue}   {RGB}{0,  70, 140}
\definecolor{thesisred}    {RGB}{180, 30,  30}
\definecolor{thesisgreen}  {RGB}{ 20,120,  60}
\definecolor{thesisorange} {RGB}{200, 90,   0}
\definecolor{thesisgray}   {RGB}{100,100, 110}
\definecolor{thesisteal}   {RGB}{  0,130, 130}
\definecolor{thesispurple} {RGB}{100,  0, 150}
\definecolor{thesisyellow} {RGB}{220,180,   0}
\definecolor{lightblue}    {RGB}{220,235,255}
\definecolor{lightgreen}   {RGB}{220,245,225}
\definecolor{lightorange}  {RGB}{255,240,220}
\definecolor{lightgray}    {RGB}{240,240,242}

%% ── pgfplots global theme ────────────────────────────────────
\pgfplotsset{
  thesis style/.style={
    width         = 0.82\textwidth,
    height        = 6.5cm,
    grid          = both,
    grid style    = {gray!15, very thin},
    tick style    = {draw=none},
    axis line style = {thesisgray},
    label style   = {font=\small},
    tick label style = {font=\scriptsize},
    legend style  = {font=\scriptsize, draw=gray!40,
                     fill=white, rounded corners=2pt},
    every axis plot/.append style = {thick},
    cycle list name = exotic,
  },
  thesis bar/.style={
    thesis style,
    ybar,
    bar width   = 8pt,
    xtick align = center,
    enlarge x limits = 0.12,
  },
}

%% ── Block-diagram node styles ────────────────────────────────
\tikzset{
  %% Generic blocks
  block/.style   = {rectangle, draw=thesisblue, thick,
                    fill=lightblue, rounded corners=3pt,
                    minimum width=2.8cm, minimum height=0.85cm,
                    align=center, font=\small},
  smallblock/.style = {rectangle, draw=thesisblue, thick,
                    fill=lightblue, rounded corners=2pt,
                    minimum width=2.0cm, minimum height=0.65cm,
                    align=center, font=\scriptsize},
  %% Process / computation blocks
  compute/.style = {rectangle, draw=thesisgreen, thick,
                    fill=lightgreen, rounded corners=3pt,
                    minimum width=2.8cm, minimum height=0.85cm,
                    align=center, font=\small},
  %% Alert / warning blocks
  alert/.style   = {rectangle, draw=thesisred, thick,
                    fill=red!6, rounded corners=3pt,
                    minimum width=2.8cm, minimum height=0.85cm,
                    align=center, font=\small},
  %% Sensor / measurement blocks
  sensor/.style  = {rectangle, draw=thesisorange, thick,
                    fill=lightorange, rounded corners=3pt,
                    minimum width=2.2cm, minimum height=0.70cm,
                    align=center, font=\scriptsize},
  %% Decision diamond
  decision/.style = {diamond, draw=thesisblue, thick,
                    fill=lightblue!60, aspect=1.8,
                    minimum width=3.0cm, minimum height=1.0cm,
                    align=center, font=\scriptsize},
  %% Summing junction
  sumjunct/.style = {circle, draw=black, thick, fill=white,
                    minimum size=0.55cm, inner sep=0pt, font=\small},
  %% Bus bar
  busbar/.style  = {draw=black, very thick, line cap=round},
  %% Arrow styles
  arr/.style     = {-{Stealth[length=2.5mm,width=1.8mm]}, thick},
  arrrev/.style  = {{Stealth[length=2.5mm,width=1.8mm]}-, thick},
  arrdbl/.style  = {{Stealth[length=2.5mm,width=1.8mm]}-
                   {Stealth[length=2.5mm,width=1.8mm]}, thick},
  arrblue/.style = {arr, thesisblue},
  arrdash/.style = {arr, dashed, thesisgray},
  %% Communication / data link
  datalink/.style = {-{Stealth[length=2.5mm,width=1.8mm]},
                    thesisgray!80, dashed, thick},
  %% Relay symbol (rectangle with protection icon label)
  relaybox/.style = {rectangle, draw=thesisred, thick,
                    fill=red!5, rounded corners=2pt,
                    minimum width=1.8cm, minimum height=0.65cm,
                    align=center, font=\scriptsize},
  %% PMU / sensor node
  pmunode/.style  = {circle, draw=thesisorange, thick,
                    fill=lightorange, minimum size=0.7cm,
                    align=center, font=\tiny},
  %% Bus node
  busnode/.style  = {circle, draw=black, thick,
                    fill=gray!15, minimum size=0.55cm,
                    inner sep=0pt, font=\scriptsize},
  %% IBR inverter box
  ibrbox/.style   = {rectangle, draw=thesisgreen, thick,
                    fill=lightgreen!70, rounded corners=2pt,
                    minimum width=1.6cm, minimum height=0.60cm,
                    align=center, font=\scriptsize},
  %% Annotation bubble
  annotbox/.style = {rectangle, draw=thesisgray!70, thin,
                    fill=white, rounded corners=2pt,
                    align=left, font=\scriptsize,
                    inner sep=3pt},
  %% Data/measurement arrow (thin dashed)
  dataarrow/.style = {-{Stealth[length=2mm,width=1.5mm]},
                      thesisgray, thin},
}

%% ── Circuit element macros ───────────────────────────────────
%% Draws a simple inductor symbol between two horizontal points
\newcommand{\drawinductor}[3][]{%
  % #1 = options, #2 = start point, #3 = end point
  \draw[#1] (#2) -- ++(0.15,0)
    \foreach \i in {1,2,3,4}{arc(90:-90:0.08) -- ++(0,0)}
    -- (#3);
}

%% ── Mathematical shorthand macros ────────────────────────────
\newcommand{\phasor}[1]{\underline{#1}}
\newcommand{\cplx}[1]{\mathbf{#1}}
\newcommand{\jj}{\mathrm{j}}
\newcommand{\dq}[1]{#1_{dq}}
\newcommand{\pu}{\ \mathrm{pu}}
\newcommand{\degs}{^{\circ}}
\newcommand{\dt}{\mathrm{d}t}
\newcommand{\Ts}{T_\mathrm{s}}

%% Sequence component notation
\newcommand{\seqpos}[1]{#1^{(+)}}
\newcommand{\seqneg}[1]{#1^{(-)}}
\newcommand{\seqzer}[1]{#1^{(0)}}
\newcommand{\Ipos}{\seqpos{I}}
\newcommand{\Ineg}{\seqneg{I}}
\newcommand{\Izer}{\seqzer{I}}
\newcommand{\Vpos}{\seqpos{V}}
\newcommand{\Vneg}{\seqneg{V}}
\newcommand{\Vzer}{\seqzer{V}}

%% Park transform notation
\newcommand{\vd}{v_d}
\newcommand{\vq}{v_q}
\newcommand{\id}{i_d}
\newcommand{\iq}{i_q}

%% Protection element notation
\newcommand{\Irated}{I_\mathrm{rated}}
\newcommand{\Ipickup}{I_\mathrm{pu}}
\newcommand{\Iscc}{I_\mathrm{sc}}
\newcommand{\ttrip}{t_\mathrm{trip}}
\newcommand{\Zreach}{Z_\mathrm{reach}}
\newcommand{\Zs}{Z_\mathrm{S}}
\newcommand{\ZR}{Z_\mathrm{R}}
\newcommand{\ZL}{Z_L}
\newcommand{\ZF}{Z_F}

%% ── Framed definition box ─────────────────────────────────────
\mdfdefinestyle{defbox}{%
  linecolor    = thesisblue,
  linewidth    = 1.2pt,
  leftmargin   = 0,
  rightmargin  = 0,
  innertopmargin    = 6pt,
  innerbottommargin = 6pt,
  innerleftmargin   = 8pt,
  innerrightmargin  = 8pt,
  backgroundcolor   = lightblue!40,
  roundcorner  = 4pt,
}

\newenvironment{thesisdef}[1]{%
  \begin{mdframed}[style=defbox]
  \textbf{Definition:} \textit{#1} ---
}{%
  \end{mdframed}
}

%% ── Framed equation highlight box ────────────────────────────
\mdfdefinestyle{eqbox}{%
  linecolor    = thesisgreen,
  linewidth    = 1pt,
  innerleftmargin  = 6pt,
  innerrightmargin = 6pt,
  backgroundcolor  = lightgreen!30,
  roundcorner  = 3pt,
}

\newenvironment{keyeq}{%
  \begin{mdframed}[style=eqbox]
}{%
  \end{mdframed}
}
